\scanpage{18}
% archive sheet number: 17
% lines 1–3
et un nb complexe \uline{non réel} $c=c'+ic''$ $(c',c''\in\unsure{R|\mathbb{R}})$ tel que $x$ et $cx\in\Gamma$. \inserted{(prendre $x\in V'\cap\Gamma$, $y\in V'\cap\Gamma$ \deleted{\unsure{si|soit} \illegible{1}} lin. ind. sur $R$)}

% lines 3
Les coordonnées de $x$ sont

% lines 4–5
\[ \alpha+\delta,\quad\beta,\quad\gamma,\quad k\delta,\qquad
\begin{aligned} &(\gamma\text{ ou }\delta\ne0\deleted{\unsure{,|)|;}}\\ &\alpha\beta\gamma\delta\in Z) \end{aligned} \]

% lines 6
celles \deleted{$\unsure{\beta|\gamma|\substack{\beta\\\gamma\\\delta}}$} de $cx$ sont

% lines 7–8
\[ \begin{aligned} &c'(\alpha+\delta)\unsure{-|+}c''\beta,\quad c''(\alpha+\delta)+c'\beta,\\
&c'\gamma-k\delta c'',\quad c''\gamma+c'k\delta \end{aligned} \]

% lines 9–14
Il faudrait donc que ce dernier vecteur puisse s'écrire $\alpha'+\delta',\beta',\gamma',k\delta',$ avec $\alpha'\beta'\gamma'\delta'\in Z$. En particulier, les premières $2$ composantes devraient être entiers, d'où on conclut que $c'$ et $c''$ seraient \uline{rationnels}.

% lines 14–16
Pour que \unsure{la|les} \deleted{deux} \deleted{deuxième} \inserted{troisième} \unsure{composante|composantes} \unsure{puisse|puissent} être \unsure{rationnelle|rationnelles}, il faut donc $\deleted{\unsure{k|\gamma}}\delta c''=0$ \deleted{\unsure{ceci|c'est} \illegible{1}}

% lines 17–18
Comme $c''\ne0$, cela implique $\delta=0$. Donc les composantes de $cx$ sont

% lines 19
\[ \alpha c'-\beta c'',\quad\beta c'+\alpha c'',\quad c'\gamma,\quad c''\gamma \]

% lines 20–23
On en tire $k\delta'=c''\gamma$, ce qui n'est possible que si les deux mbres sont nuls, d'où $\gamma=0$. Mais, \deleted{\unsure{donc}} \inserted{comme $\gamma=\delta=0$} $x$ serait dans \unsure{l'espace|l'esp.} engendré par $\unsure{e_0|e_1}$ et $ie_1$, absurde
